% Part: second-order-logic
% Chapter: sol-and-sets
% Section: introduction

\documentclass[../../../include/open-logic-section]{subfiles}

\begin{document}

\olfileid[fr]{sol}{set}{int}

\olsection{Introduction}

La logique du second ordre permet de quantifier
sur les sous-ensembles du !!{domain} ainsi que
sur les fonctions. Il est donc naturel qu'elle
permette d'exprimer au moins une partie de la
théorie des ensembles. Par « exprimer », nous
entendons formuler des propriétés et des
énoncés ensemblistes en logique du second ordre,
sans introduire de vocabulaire non logique
propre aux ensembles, tel que le !!{predicate}
d'appartenance de la théorie des ensembles.
Nous pouvons, par exemple, définir l'union et
l'intersection d'ensembles, ainsi que la relation
d'inclusion; nous pouvons aussi comparer leurs
tailles et énoncer des résultats comme le
théorème de Cantor.

\end{document}
